\chapter{Conclusion}

\section{Conclusion and Main Contribution}
This thesis presented Intent-Spawner, a bounded pre-spawn decision-support system for selecting JupyterHub environments from an administrator-curated catalog. P2 is the main proposed method. It validates a structured representation of current workload intent, uses BM25 and dense retrieval to generate candidates, fuses ranks with RRF, applies hard constraints before preference ranking, resolves selected identifiers from trusted catalog data, and binds the previewed choice to an authenticated user's confirmation. This design preserves a clear policy boundary and keeps the platform's ordinary admission and scheduling mechanisms authoritative.

The work also specified Protocol-v5 as a set of distinct experiment streams. That separation matters because recommendation relevance, robustness to prompt variation, human experience, resource adequacy, functional image behavior, and storage accounting are not interchangeable outcomes. The current repository contains development and local-environment observations across several of these streams, but the consolidated analysis is incomplete. Human participants have not been observed, several resource and image claims require audit and context-specific interpretation, and confirmatory decisions remain \texttt{NOT\_EXECUTED}. Therefore, the principal empirical conclusion is a statement about evidence status: the prototype and evaluation infrastructure exist, but the registered claims have not yet been confirmed.

The system integration is the thesis's primary contribution. It shows how recommendation techniques can be applied at a practical JupyterHub configuration boundary while constraining effects through trusted catalog resolution, deterministic feasibility rules, preview, and user confirmation. It does not show that intent predicts actual peak demand, that P2 is statistically superior to P1, or that P2 improves participant outcomes over B0. The local E4/E5 observations are useful for refining protocols and exposing operational questions, but they do not justify a broad production claim.

\section{Limitations and Practical Implications}
Several limitations remain. First, offline recommendation evidence is development-stage and the current analysis package has metric and provenance reconciliation work outstanding. Second, the independent semantic sample is based on workload families, and 18 development families cannot support broad conclusions across all notebook domains. Third, E3 has no participant data, so the value of recommendations to real users is unknown. Fourth, E4's OrbStack environment differs from a representative multi-node cluster and must be reviewed for measurement-source, plan, and execution scope. Fifth, E5 tests only a small approved catalog, a particular platform, and bounded functional probes. Manifest layer bytes do not equal node disk use. Finally, a prototype preview-state design needs shared durable storage before it can be treated as suitable for multi-replica production deployment.

The feature scope is also intentionally narrow. P2 does not inspect live cluster capacity, schedule a Pod, guarantee successful admission, or replace runtime resource controllers. The dynamic-sizing component remains bounded prototype work. The reprovisioning workflow stops and recreates a server and cannot preserve in-memory execution state. P3 is not retained. These boundaries should remain visible in any abstract, defense presentation, or future system description.

For an administrator, the prototype suggests a workflow in which a candidate catalog is treated as versioned operational data rather than as an informal list of image names. Capability metadata should be testable, image references should be immutable, and policy should be evaluated when a preview is created and again when it is applied. An administrator considering deployment should also decide how unsupported or ambiguous requests are communicated and whether any manual override remains available under local policy.

For a user, the intended value is less translation work. A user can express the task in familiar language and inspect a suggestion together with resource and capability information. The user should still be able to spot an unsuitable candidate, revise the intent, and understand when no approved profile satisfies a requirement. This interaction can reduce friction only if recommendations are sufficiently relevant and the preview is understandable. Since human evaluation has not been completed, these statements describe design intent, not demonstrated user benefit.

For a research evaluator, the project illustrates the importance of separating algorithmic evidence from integration and operational evidence. A recommender can be implemented correctly and still have weak rankings. A relevant candidate can fail a functional probe. A working Pod can still consume more resources than necessary. A successful task can still be harder for users to complete. The evaluation streams therefore provide a more useful roadmap than a single all-purpose score.

\section{Prioritized Future Work}
The next work should address evidence integrity before expanding the method. The first priority is to resolve the Protocol-v5 analysis package: reconcile E1 and E5 metric definitions, trace every derived value to raw observations, review dirty-tree provenance, and produce an eligible evidence-selection report. The second priority is to run frozen E1/E2 confirmation on sealed family-level cases. No implementation or configuration tuning should use those final cases.

The third priority is E3. The study should complete ethics and consent preparation, recruit an appropriate participant sample, counterbalance conditions, preserve privacy, and analyze participant-level outcomes using the registered plan. The fourth priority is E4: close the plan-versus-observation audit, declare the eligible cgroup/Kubernetes resource source and conditions, and execute a protocol-conformant study on the intended cluster class. The fifth priority is E5 validation at the intended platform and catalog scope. Where eight- or sixteen-image catalogs do not exist, those cells should remain explicitly \texttt{NOT\_EXECUTED} until suitable approved images are available.

Only after these work packages should development focus on improving top-one ranking, expanding catalog metadata, or studying optional dynamic sizing. P2 changes should be evaluated on development data and refrozen before any new confirmation. Multi-node, multi-tenant evaluation should examine quota contention, registry pulls, scheduling variability, policy drift, and failure recovery. Productionization should move preview state to a shared durable store and test atomic token consumption across replicas. Any future use of behavioral history requires a separate consent, retention, user-splitting, and privacy threat model.

Intent-Spawner is best understood as an auditable, human-confirmed recommender for a bounded set of notebook environments. Its strongest demonstrated result is the concrete integration of intent parsing, retrieval, constraints, trusted selection, and pre-spawn application, together with an evaluation protocol that makes evidence boundaries explicit. The empirical record currently supports descriptive development and local observations only. Confirmatory comparisons and general deployment claims remain open work. Maintaining this distinction is necessary for an accurate final thesis and provides a clear path from the current prototype to defensible future conclusions.
